\ExamPart{Câu tự luận}{3}{ }

\NQuestion Giải bất phương trình $\left(x^2-5x+6\right)\left(x^2-4\right)\le 0$.
\AnswerLine{$[-2;3]$}
\SolutionBlock{Ta phân tích
\[
\left(x^2-5x+6\right)\left(x^2-4\right)=(x-2)(x-3)(x-2)(x+2)=(x-2)^2(x-3)(x+2).
\]
Vì $(x-2)^2\ge0$ với mọi $x$, nên dấu của biểu thức phụ thuộc vào $(x-3)(x+2)$, đồng thời $x=2$ cũng làm biểu thức bằng $0$. Ta có
\[
(x-3)(x+2)\le0 \Leftrightarrow x\in[-2;3].
\]
Do đó tập nghiệm của bất phương trình đã cho là
\[
[-2;3].
\]}

\NQuestion Tìm tất cả các giá trị của $m$ để phương trình
\[
x^2-2(m-1)x+m^2-4m+3=0
\]
có hai nghiệm phân biệt $x_1,x_2$ thỏa mãn $x_1+x_2<4$.
\AnswerLine{$1<m<3$}
\SolutionBlock{Theo hệ thức Viète,
\[
x_1+x_2=2(m-1).
\]
Điều kiện $x_1+x_2<4$ cho
\[
2(m-1)<4\Leftrightarrow m<3.
\]
Để phương trình có hai nghiệm phân biệt, cần và đủ $\Delta>0$. Ta có
\[
\Delta=[-2(m-1)]^2-4(m^2-4m+3)=4(m-1)^2-4(m^2-4m+3)=8m-8.
\]
Suy ra
\[
\Delta>0\Leftrightarrow 8m-8>0\Leftrightarrow m>1.
\]
Kết hợp lại, ta được
\[
1<m<3.
\]}

\NQuestion Trong mặt phẳng tọa độ $Oxy$, cho ba điểm $A(1;1),\ B(5;3),\ C(3;-3)$.
\begin{SubQuestions}
  \item Tính tọa độ các vectơ $\overrightarrow{AB}$ và $\overrightarrow{AC}$.
  \item Chứng minh rằng tam giác $ABC$ vuông tại $A$.
  \item Viết phương trình đường thẳng $BC$.
  \item Tính khoảng cách từ điểm $A$ đến đường thẳng $BC$.
\end{SubQuestions}
\AnswerLine{$\overrightarrow{AB}=(4;2),\ \overrightarrow{AC}=(2;-4),\ BC:3x-y-12=0,\ d(A,BC)=\sqrt{10}$}
\SolutionBlock{\begin{SubQuestions}
\item Ta có
\[
\overrightarrow{AB}=(5-1;3-1)=(4;2),\qquad \overrightarrow{AC}=(3-1;-3-1)=(2;-4).
\]
\item Tính tích vô hướng:
\[
\overrightarrow{AB}\cdot\overrightarrow{AC}=4\cdot2+2\cdot(-4)=8-8=0.
\]
Vì tích vô hướng bằng $0$ nên $\overrightarrow{AB}\perp\overrightarrow{AC}$. Vậy tam giác $ABC$ vuông tại $A$.
\item Đường thẳng $BC$ đi qua $B(5;3)$ và $C(3;-3)$ nên có hệ số góc
\[
m=\dfrac{-3-3}{3-5}=3.
\]
Phương trình qua $B(5;3)$ là
\[
y-3=3(x-5)\Leftrightarrow 3x-y-12=0.
\]
\item Khoảng cách từ $A(1;1)$ đến $BC:3x-y-12=0$ là
\[
d(A,BC)=\dfrac{|3\cdot1-1-12|}{\sqrt{3^2+(-1)^2}}=\dfrac{10}{\sqrt{10}}=\sqrt{10}.
\]
\end{SubQuestions}}

\NQuestion Trong mặt phẳng tọa độ $Oxy$, cho điểm $M(2;-1)$ và đường thẳng
\[
d:x-2y+3=0.
\]
\begin{SubQuestions}
  \item Viết phương trình đường thẳng $\Delta$ đi qua $M$ và song song với $d$.
  \item Viết phương trình đường thẳng $\Delta'$ đi qua $M$ và vuông góc với $d$.
  \item Tính khoảng cách từ $M$ đến đường thẳng $d$.
\end{SubQuestions}
\AnswerLine{$\Delta:x-2y-4=0,\ \Delta':2x+y-3=0,\ d(M,d)=\dfrac{7}{\sqrt5}$}
\SolutionBlock{\begin{SubQuestions}
\item Vì $\Delta\parallel d$ nên $\Delta$ có cùng vectơ pháp tuyến với $d$, do đó có dạng
\[
x-2y+c=0.
\]
Thay $M(2;-1)$ vào, ta được $2-2(-1)+c=0\Rightarrow c=-4$. Vậy
\[
\Delta:x-2y-4=0.
\]
\item Đường thẳng $d:x-2y+3=0$ có vectơ pháp tuyến là $(1;-2)$, nên có thể lấy một vectơ chỉ phương là $(2;1)$. Vì $\Delta'\perp d$ nên $\Delta'$ có vectơ pháp tuyến là $(2;1)$, do đó có dạng
\[
2x+y+c=0.
\]
Thay $M(2;-1)$ vào, ta được $2\cdot2+(-1)+c=0\Rightarrow c=-3$. Vậy
\[
\Delta':2x+y-3=0.
\]
\item Khoảng cách từ $M$ đến $d$ là
\[
d(M,d)=\dfrac{|2-2(-1)+3|}{\sqrt{1^2+(-2)^2}}=\dfrac{7}{\sqrt5}.
\]
\end{SubQuestions}}
